\documentclass[11pt]{article}

\usepackage{../ppackage}
\usepackage{bbm}
\usepackage{import}
\usepackage{refstyle}





\usepackage{tikz}
\usetikzlibrary{arrows.meta,patterns}

\usepackage{../../../Notes/tikzit}
\usepackage{../../../Notes/utf8math}

\input{../../../Notes/TIKZ/digraph.tikzstyles}



\usepackage{url}


\newcommand{\solution}{
\bigskip\noindent
	\textbf{Solution: \\}}
	


\DeclareMathOperator{\size}{size}
\DeclareMathOperator{\conv}{conv}
\newcommand{\SV}{\mathrm{SV}}
\newcommand{\bigO}{O}
\newcommand{\cut}{\mathrm{cut}}
\newcommand{\LLL}{\mathrm{LLL}}
\newcommand{\setR}{\mathbb{R}}
\newcommand{\setZ}{\mathbb{Z}}
\newcommand{\setQ}{\mathbb{Q}}
\newcommand{\setC}{\mathbb{C}}
\newcommand{\setN}{\mathbb{N}}
\newcommand{\wt}[1]{\widetilde{#1}}
\newcommand{\opt}{{\sc 0/1-opt}\xspace}
\newcommand{\aug}{{\sc 0/1-aug}\xspace}
\newcommand{\psep}{{\sc 0/1-psep}\xspace}
\newcommand{\sep}{{\sc 0/1-sep}\xspace}
\newcommand{\fopt}{{\sc 0/1-testopt\xspace} }

\newcommand{\hpp}{\mathrm{HPP}}
\newcommand{\nodes}{\mathcal{V}}
\newcommand{\vol}{\mathrm{vol}}
\newcommand{\diag}{\mathrm{diag}}
\newcommand{\arcs}{\mathcal{A}}
\newcommand{\edges}{\mathcal{E}}
\newcommand{\paths}{\mathscr{P}}
\newcommand{\cycles}{\mathcal{C}}




\newcommand{\K}{{\mathcal K}}
\newcommand{\A}{{A}}
\newcommand{\B}{{B}}
\newcommand{\T}{\mathscr{T}}
\newcommand{\eE}{\mathscr{E}}
\newcommand{\eS}{\mathscr{S}}
\newcommand{\eP}{\mathscr{P}}
\newcommand{\eM}{\mathscr{M}}



\newcommand{\transp}{^{\mathrm{T}}}

\newcommand{\smallmat}[1]{\left( \begin{smallmatrix} #1 \end{smallmatrix}\right)}

\newcommand{\mat}[1]{ \begin{pmatrix} #1 \end{pmatrix}}
\newcommand{\smat}[1]{ \big(\begin{smallmatrix} #1 \end{smallmatrix}\big)}

\newcommand{\pc}{\mathscr{P}}
\newcommand{\ob}{\mathscr{O}}
\newcommand{\odds}{\mathscr{W}}
\newcommand{\up}{\mathscr{U}}
\newcommand{\ef}{\mathscr{F}}
\newcommand{\eh}{\mathscr{H}}
\newcommand{\ev}{\mathscr{V}}
\newcommand{\ec}{\mathscr{C}}
\newcommand{\eu}{\mathscr{U}}

\newcommand{\lex}{\mathrm{lex}}

\renewcommand{\leq}{\leqslant}
\renewcommand{\geq}{\geqslant}









\newcommand{\linhull}{\mathrm{lin.hull}}
\newcommand{\affhull}{\mathrm{affine.hull}}
\newcommand{\charcone}{\mathrm{char.cone}}
\newcommand{\cone}{\mathrm{cone}}
\newcommand{\rank}{\mathrm{rank}}
\newcommand{\wb}[1]{\overline{#1}}



\usepackage{enumerate}

      
\institute{\'Ecole Polytechnique F\'ed\'erale de Lausanne}
\lecture{Discrete Optimization}
\faculty{Prof. Friedrich Eisenbrand}
\term{Spring 2026}
\publishdate{May 19, 2026}
\duedate{ }
\problemset{Problem Set~12}

\begin{document}
\makeheader

\begin{enumerate}[1)]

\item \label{Ex:Node-Pot} Let $G = (V,A)$ be a directed graph and $c: A ⟶ \mathbb{R}$ be a length function. Consider the following set of linear inequalities with variables $π$:
  \begin{equation}\label{eq:pot5}
    ∀ (u,v) ∈ A: π(v) ≤ π(u) + c_{uv}. 
  \end{equation}
  \begin{enumerate}[i)]
  \item   Show that this system of inequalities has a feasible solution if and only if $G$ does not possess a cycle of negative weight.
  \item  Interpret a cycle of negative weight as a Farkas certificate of infeasibility.
  \item Let $s ∈ V$ be a vertex from where each other vertex of the graph can be reached.  Consider the linear program which maximizes $∑_{v ∈ V} π(v) $ subject to \ref{eq:pot5} and the additional constraint $π(s) = 0$. Show that the linear program has a unique optimal solution $π^*$. 
  
  \emph{Hint: What is $π^*$ in terms of distances from $s$? }
\end{enumerate}

\item \label{ex:count:paths} Let $G = (V,A)$ be a directed graph. The \emph{directed adjacency matrix} $A ∈ \{0,1\}^{V \times V}$ is  defined as 
\[
a_{uv} =
\begin{cases}
1, & \text{if there is an arc } u\to v,\\
0, & \text{otherwise.}
\end{cases}
\]
So \(A\) encodes exactly which directed edges, or arcs, are present in the graph.
\begin{enumerate}[i)] 
\item Show by induction that for  $k ∈ ℕ_+$ that  $A^k$ encodes the number of directed walks of length $k$ from $u$ to $v$. In other words
  \[
a^k_{uv} =  \text{ number of directed walks of length $k$ from $u$ to $v$.} 
\]
\item Show that $A^k$ can be computed with $O(\log k ⋅ n^{2.805})$ arithmetic operations.
\end{enumerate}

\item \label{triangle-detection} Let $G = (V,E)$ be an undirected graph. A \emph{triangle} in $G$ is a set of distinct vertices $u,v,w ∈ V$ that are pairwise connected by an edge.
  Design a randomized algorithm that decides whether $G$ contains a triangle. You can assume that $G$ is represented by its \emph{adjacency matrix} $A ∈ \{0,1\}^{V \times V}$ where  $a_{uv} = 1$ if $uv ∈E$ and otherwise $a_{uv} = 0$.  
  It should satisfy the following.
  \begin{enumerate}[i)] 
  \item If $G$ does not contain a triangle, then it always asserts \emph{No}.
  \item If the algorithm asserts \emph{Yes}, then $G$ contains a triangle.
  \item The probability of the assertion \emph{No} provided that $G$ has a triangle should be bounded by $1/2$.
  \item The algorithm should perform $O(n^{2.805})$ arithmetic operations on integers of size $O(\log \log n)$. 
  \end{enumerate}

\item The \emph{node-edge incidence matrix} $A_G ∈ \{0, \pm 1\}^{V \times A}$ of a directed graph $G = (V,A)$   is defined as
  \begin{displaymath}
    a_{ua} = \begin{cases}
      1, & \text{if  } a= (x,u) \text{ for some } x ∈V\\
      -1, & \text{if  } a= (u,x) \text{ for some } x ∈V\\
0, & \text{otherwise.}
    \end{cases}    
  \end{displaymath}
  Show that $A_G$ is totally unimodular. 
\end{enumerate}





  

\end{document}

%%% Local Variables:
%%% mode: latex
%%% TeX-master: t
%%% End:
